\section{A scalar-curvature identity for graph stabilization}
\label{sec:graph-calculus}

One shortening step increases a metric in the direction of a single
function's differential and adds one circle. The circle warp compensates
for the scalar terms introduced by the graph. Warped-circle formulas are
also central to the stabilization used in minimal slicing
\cite[Lemma~2.5]{SchoenYau2017}; here we compute their combination with a
rank-one graph change before choosing the graph by a nonlinear equation.

We use the Laplacian $\Delta=\diver\nabla$ and the convention for which
round spheres have positive scalar curvature. The base $(B,g)$ may have
any finite dimension and may already contain warped circle factors. All
derivatives not explicitly assigned another metric are taken with respect
to $g$.

Let $f\in C^\infty(B)$ and let $N:[0,\infty)\to(0,\infty)$ be smooth and constant near zero. Set
\begin{equation}\label{eq:graph-quantities}
 \begin{gathered}
 t=\abs{df}_g,\qquad W=(1+N(t)^2t^2)^{1/2},\qquad v=W^{-1},\\
 U=vN(t)\nabla f,\qquad r=\abs{U}_g=vN(t)t,\qquad
 D=-\diver_g U.
 \end{gathered}
\end{equation}
Although $t$ and $r$ need not be smooth where $df=0$, the functions $N(t)$
and $v$ and the vector field $U$ are smooth there. Indeed, $N(t)$ is
constant near such points, and the other two are smooth functions of $df$.
We use identities involving $\log t$ or $\log r$ only where $t>0$; the
resulting scalar estimates extend by continuity. Consider
\begin{equation}\label{eq:graph-metric-update}
 g^+=g+N(t)^2df^2+v^2d\theta^2
 \quad\text{on }B\times\T.
\end{equation}
The first two terms form the metric induced on the graph of $f$ in
$(B\times\R,g+N(t)^2dz^2)$, where $N(t)$ is regarded as a fixed function of the base point, independent of $z$.

\begin{proposition}[Exact scalar-curvature identity]\label{prop:scalar-update}
At a point where $t>0$, choose a $g$-orthonormal frame
$e_1=\nabla f/t,e_2,\ldots,e_{\dim B}$, and put
\[
 \mathfrak p=\frac{tN'(t)}{N(t)},\qquad
 A=d\log t,\qquad
 C_{jk}=t^{-1}\Hess_g f(e_j,e_k)\quad(j,k>1).
\]
Then
\begin{equation}\label{eq:scalar-exact}
 \begin{split}
 \Scal_{g^+}-\Scal_g
 ={}&D^2-2U(D)+r^2\abs{C}^2+(e_1r)^2\\
 &+2r^2(1+\mathfrak p)(v^2-\mathfrak p r^2)
       \sum_{j>1}A(e_j)^2.
 \end{split}
\end{equation}
In particular, the final term is nonnegative whenever
$\mathfrak p\ge0$ and $\mathfrak p N(t)^2t^2\le1$.
For $N\equiv1$, one also has, everywhere,
\begin{equation}\label{eq:ordinary-scalar-update}
 \Scal_{g^+}-\Scal_g=\abs{II}^2+D^2-2U(D),
\end{equation}
where $II$ is the second fundamental form of the ordinary product graph.
The inverse graph metric on base covectors and the volume form of $g^+$ are
\begin{equation}\label{eq:graph-inverse-volume}
 P=g^{-1}-U\otimes U,
 \qquad dV_{g^+}=dV_g\,d\theta.
\end{equation}
\end{proposition}

\begin{proof}
Write $E=N^{-1}\partial_z$ and give the graph the upward unit normal
$\nu=-U+vE$. Our convention for the second fundamental form is
$II(X,Y)=\ip{\nabla_X\nu,Y}$, so its trace is
\begin{equation}\label{eq:graph-mean-curvature}
 H=D-U(\log N).
\end{equation}
The vertical Killing field $\partial_z$ has normal speed $Nv$ along the
graph, and its tangential part projects to $NvU$ on the base. Vertical
translation preserves the graph's mean curvature as a function of that
base point. The normal and tangential variation terms must therefore
cancel, giving
\begin{equation}\label{eq:graph-jacobi}
 \bigl(\Delta_{\rm gr}+\abs{II}^2+\Ric_{\rm amb}(\nu,\nu)\bigr)(Nv)
       =Nv\,U(H).
\end{equation}
Here a normal variation of speed $\psi$ changes $H$ by the negative of the
Jacobi operator applied to $\psi$, while a tangential variation
differentiates $H$ in its tangential direction. This fixes the sign on the
right of \eqref{eq:graph-jacobi}.

The classical warped-product formulas \cite[Section~7]{BishopONeill1969},
together with the Gauss equation, give, respectively,
\[
 \begin{aligned}
 \Scal_{\rm amb}&=\Scal_g-2\Delta_gN/N,\\
 \Scal_{\rm gr}&=\Scal_{\rm amb}-2\Ric_{\rm amb}(\nu,\nu)
                         +H^2-\abs{II}^2,\\
 \Scal_{g^+}&=\Scal_{\rm gr}-2\Delta_{\rm gr}v/v.
 \end{aligned}
\]
Expanding $\Delta_{\rm gr}(Nv)$ in \eqref{eq:graph-jacobi} therefore yields
\begin{equation}\label{eq:scalar-before-cancellation}
 \begin{split}
 \Scal_{g^+}-\Scal_g
 ={}&\abs{II}^2+H^2-2U(H)\\
 &+2\frac{(\Delta_{\rm gr}-\Delta_g)N}{N}
   +4\ip{d\log N,d\log v}_{P}.
 \end{split}
\end{equation}
The formula for $P$ in \eqref{eq:graph-inverse-volume} follows directly by inverting the rank-one addition $g+N^2df^2$.

Put $b_0=d\log N$. For a function on the base,
$\Delta_{\rm gr}N$ is its ambient Hessian traced over the graph, minus $H\nu(N)$. Since
$\Hess_{\rm amb}N(E,E)=N\abs{b_0}^2$, this gives
\begin{equation}\label{eq:graph-laplacian-correction}
 \frac{(\Delta_{\rm gr}-\Delta_g)N}{N}
 =-\Hess_g\log N(U,U)-b_0(U)^2
       +r^2\abs{b_0}^2+H b_0(U).
\end{equation}
Substitute $H=D-b_0(U)$ in \eqref{eq:scalar-before-cancellation}. The Hessian terms cancel because
\[
 U\bigl(b_0(U)\bigr)
   =\Hess_g\log N(U,U)+b_0(\nabla^g_UU).
\]
The result is
\begin{equation}\label{eq:scalar-intermediate}
 \begin{split}
 \Scal_{g^+}-\Scal_g
 ={}&\abs{II}^2+D^2-2U(D)
       +2b_0(\nabla^g_UU)-3b_0(U)^2\\
 &+2r^2\abs{b_0}^2+4\ip{b_0,d\log v}_{P}.
 \end{split}
\end{equation}
When $N\equiv1$, all terms involving $b_0$ vanish. This proves
\eqref{eq:ordinary-scalar-update} without any condition on $D$ other than its definition.

For the remaining calculation use the graph-orthonormal frame
$T=ve_1+rE,e_2,\ldots,e_{\dim B}$. For base vector fields lifted independently of $z$, the ambient connection satisfies
\[
 \nabla_X^{\rm amb}E=0,\qquad
 \nabla_E^{\rm amb}X=b_0(X)E,\qquad
 \nabla_E^{\rm amb}E=-b_0^{\sharp}.
\]
Differentiating $\nu=-U+vE$ gives
\begin{equation}\label{eq:graph-second-form}
 \begin{aligned}
 II(T,T)&=-r\bigl[v^2(1+\mathfrak p)+\mathfrak p\bigr]A(e_1),\\
 II(T,e_j)&=-rv(1+\mathfrak p)A(e_j),\qquad j>1,\\
 II(e_j,e_k)&=-rC_{jk},\qquad j,k>1.
 \end{aligned}
\end{equation}
Moreover,
\begin{equation}\label{eq:graph-log-derivatives}
 \begin{gathered}
 b_0=\mathfrak p A,\qquad
 d\log r=v^2(1+\mathfrak p)A,\qquad
 d\log v=-r^2(1+\mathfrak p)A,\\
 \nabla^g_UU=r^2\left(v^2(1+\mathfrak p)A(e_1)e_1
                         +\sum_{j>1}A(e_j)e_j\right).
 \end{gathered}
\end{equation}
For the last equality, symmetry of $\Hess f$ gives
$\nabla_{e_1}e_1=\sum_{j>1}A(e_j)e_j$.

After substitution in \eqref{eq:scalar-intermediate}, the $C$ terms give
$r^2\abs{C}^2$. The remaining terms are squares of the component of $A$
along $e_1$ and of its components orthogonal to $e_1$. After division by
$r^2$, the coefficient of $A(e_1)^2$ is
\[
 \begin{split}
 &\bigl[v^2(1+\mathfrak p)+\mathfrak p\bigr]^2
       -2\mathfrak p v^2(1+\mathfrak p)-\mathfrak p^2\\
 &\hspace{35mm}=v^4(1+\mathfrak p)^2.
 \end{split}
\]
For every $j>1$, the corresponding coefficient is
\[
 2v^2(1+\mathfrak p)^2-2\mathfrak p(1+\mathfrak p)
       =2(1+\mathfrak p)(v^2-\mathfrak p r^2).
\]
Since $e_1r=rv^2(1+\mathfrak p)A(e_1)$, these are precisely
\eqref{eq:scalar-exact}. Also
\[
 v^2-\mathfrak p r^2=v^2\bigl(1-\mathfrak p N(t)^2t^2\bigr),
\]
which proves the positivity assertion. Finally,
$\det(g+N^2df^2)=W^2\det g$ and $vW=1$ prove the volume identity.
\end{proof}

The inverse-metric decrease will later let us sum the costs of many graph
changes in a fixed $n$-dimensional space. First we must choose each graph
so that the remaining term $D^2-2U(D)$ can be bounded from below.
The next section proves existence for the equation used for this choice.
